%%%%%%%%%%%%%%%%%%%%%%%%%%%%%%%%%%%%%%%%%%%%%%%%%%%%%%%%%%%%%%%%%%%%%%%%%%%%%%%%
%%% TIKZLIBRARY stap.components %%%%%%%%%%%%%%%%%%%%%%%%%%%%%%%%%%%%%%%%%%%%%%%%%
%   Building blocks for drawing arithmetization-oriented / ZK-friendly          %
%   permutations (S-boxes, linear layers, operators, wires, braces).            %
%   Load with:  \usetikzlibrary{stap.components}                                %
%                                                                              %
%   Convention: operators are ALL-CAPS node styles; boxes, state, wires and    %
%   spacers are lowercase. Figures build their own concrete styles on top of    %
%   the base elements here (e.g. roundbox = vbox={1.2cm, fill=typeone}).        %
%%%%%%%%%%%%%%%%%%%%%%%%%%%%%%%%%%%%%%%%%%%%%%%%%%%%%%%%%%%%%%%%%%%%%%%%%%%%%%%%

% fit           -> layerbox auto-sizing;   calc  -> $(a)!.5!(b)$ in braces/figures
% arrows.meta   -> wire arrow tip;         positioning -> right=of / below=of
% decorations.pathreplacing -> \hbrace / \vbrace
\usetikzlibrary{fit, calc, arrows.meta, positioning, decorations.pathreplacing}

%%% ================================================================
%%% OPERATORS  (small inline operation nodes)
%%% ================================================================

\tikzset{
  %% shared base for all operator symbols -- inherit, don't use directly
  op base/.style={draw, minimum size=0.35cm, inner sep=0pt},
  %% NOOP: invisible placeholder of operator size (keeps wires aligned)
  NOOP/.style={op base, draw=none, fill=none,},
  %% XOR: circle with a cross
  XOR/.style={op base, circle,
    append after command={
      [shorten >=\pgflinewidth, shorten <=\pgflinewidth]
      (\tikzlastnode.north) edge (\tikzlastnode.south)
      (\tikzlastnode.east)  edge (\tikzlastnode.west)}},
  %% AND: circle with an x (boolean AND)
  AND/.style={op base, circle,
    append after command={
      [shorten >=\pgflinewidth, shorten <=\pgflinewidth]
      (\tikzlastnode.north west) edge (\tikzlastnode.south east)
      (\tikzlastnode.north east) edge (\tikzlastnode.south west)}},
  %% PLUS: square with + (any field/modular addition, e.g. round constants)
  PLUS/.style={op base, rectangle,
    append after command={
      [shorten >=\pgflinewidth, shorten <=\pgflinewidth]
      (\tikzlastnode.north) edge (\tikzlastnode.south)
      (\tikzlastnode.east)  edge (\tikzlastnode.west)}},
  %% MINUS: square with -
  MINUS/.style={op base, rectangle,
    append after command={
      [shorten >=\pgflinewidth, shorten <=\pgflinewidth]
      (\tikzlastnode.east) edge (\tikzlastnode.west)}},
  %% MUL: square with an x (field multiplication)
  MUL/.style={op base, rectangle,
    append after command={
      [shorten >=\pgflinewidth, shorten <=\pgflinewidth]
      (\tikzlastnode.north west) edge (\tikzlastnode.south east)
      (\tikzlastnode.north east) edge (\tikzlastnode.south west)}},
}


%%% ================================================================
%%% BOXES  (labeled function boxes: S-box, F, pi, linear layers, ...)
%%% ================================================================

\tikzset{
  %% base box -- inherit this, don't use directly
  box base/.style={draw, inner sep=0.5pt, text height=2ex, align=center, fill=gray!10},
  %% bbox={width}{height}: box of an explicit size
  bbox/.style 2 args={box base, minimum width=#1, minimum height=#2},
  %% sbox: square box for an S-box / single-wire function / constant
  sbox/.style={bbox={0.8cm}{0.8cm}},
  %% vbox={height}: fixed-width, custom-height box (e.g. round-function box)
  vbox/.style = {bbox={0.7cm}{#1}},
  %% hbox={width}: fixed-height, custom-width box
  hbox/.style = {bbox={#1}{0.7cm}},
  %% layerbox={(a)(b)...}{label}: box auto-sized to the bounding box of the
  %% listed nodes -- the idiom for a layer spanning a whole state.
  %% usage: \node[layerbox={(l-1)(l-t)}{$\M$}, below=of ...] (name) {};
  layerbox/.style 2 args={box base, fit={#1}, label=center:{#2}},
}


%%% ================================================================
%%% SPACERS / LABELS  (invisible nodes for alignment and dots)
%%% ================================================================

\tikzset{
  %% empty: spacer / label node (e.g. holding \dots or \vdots in a state)
  empty/.style      = {align=center, minimum size=1cm},
  %% emptysmall: a tighter spacer
  emptysmall/.style = {align=center, minimum size=1mm},
  %% note: small gray caption text (e.g. under a brace)
  note/.style       = {color=gray, font=\scriptsize},
}


%%% ================================================================
%%% STATE WIRES
%%% ================================================================

\tikzset{
  %% state: a state-element node (input/output value on a wire)
  state/.style = {minimum size=0.5cm, inner sep=0pt, outer sep=0pt},
  %% wire: directed connection with a small Latex arrow tip
  wire/.style        = {-{Latex[length=5pt]},},
  wire dashed/.style = {wire, dashed},
}


%%% ================================================================
%%% BRACES
%%% ================================================================

%% \hbrace[left=, right=, raise=]{P1}{P2}{P3}
%%   horizontal brace at y(P3), from x(P1)-left to x(P2)+right
\pgfkeys{/hbrace/.is family, /hbrace,
  left/.initial=0pt, right/.initial=0pt, raise/.initial=4pt,
  color/.initial=gray, font/.initial=\scriptsize}
\newcommand{\hbrace}[4][]{%
  \pgfkeys{/hbrace, left=0pt, right=0pt, raise=4pt, color=gray, font=\scriptsize, #1}%
  \draw[decorate, decoration={brace, mirror},
        color=\pgfkeysvalueof{/hbrace/color}, font=\pgfkeysvalueof{/hbrace/font}]
    ($(#2 |- #4)+(-\pgfkeysvalueof{/hbrace/left},-\pgfkeysvalueof{/hbrace/raise})$) --
    ($(#3 |- #4)+(\pgfkeysvalueof{/hbrace/right},-\pgfkeysvalueof{/hbrace/raise})$)%
}

%% \vbrace[above=, below=, raise=]{P1}{P2}{P3}
%%   vertical brace to the right of x(P3)+raise, from y(P2)-below upward to y(P1)+above
%%   P1 = top node, P2 = bottom node, P3 = x reference (typically rightmost node)
%%   brace opens leftward toward content; node[midway, right=..] places label outside
\pgfkeys{/vbrace/.is family, /vbrace,
  above/.initial=0pt, below/.initial=0pt, raise/.initial=4pt,
  color/.initial=gray, font/.initial=\scriptsize}
\newcommand{\vbrace}[4][]{%
  \pgfkeys{/vbrace, above=0pt, below=0pt, raise=4pt, color=gray, font=\scriptsize, #1}%
  \draw[decorate, decoration={brace},
        color=\pgfkeysvalueof{/vbrace/color}, font=\pgfkeysvalueof{/vbrace/font}]
    ($(#3 -| #4)+(+\pgfkeysvalueof{/vbrace/raise},-\pgfkeysvalueof{/vbrace/below})$) --
    ($(#2 -| #4)+(+\pgfkeysvalueof{/vbrace/raise},+\pgfkeysvalueof{/vbrace/above})$)%
}
